\subsection{Riscos e contingência}

O projeto envolve desafios teóricos, algorítmicos e 
experimentais, de modo que é necessário antecipar possíveis 
dificuldades e definir estratégias de resposta. 
A seguir, são apresentados os principais riscos e as 
contingências previstas.

\subsubsection{Formalização teórica da operação de split}

Um risco central é a dificuldade de formalizar a divisão 
de um nó da \textit{max tree} em múltiplos nós-filhos, 
especialmente no que tange à redistribuição de 
componentes conexas e de arestas após uma alteração 
local na imagem. Essa dificuldade pode comprometer 
tanto a correção do modelo quanto a clareza da 
demonstração matemática.

Como contingência, a formalização será conduzida em 
etapas progressivas: primeiro, com transformações simples 
e casos particulares; em seguida, com generalização para 
situações mais complexas. A construção de contraexemplos 
e casos limite será usada para validar a consistência 
do arcabouço teórico.

\subsubsection{Eficiência e corretude do algoritmo}

O operador de \textit{split} pode não atingir o nível de 
eficiência desejado, ou ainda pode exigir estruturas de 
dados e procedimentos de manutenção mais complexos do 
que o inicialmente previsto. Nesse caso, o algoritmo 
pode perder vantagem em relação à reconstrução 
completa da árvore.

Como contingência, será adotada uma estratégia 
incremental de desenvolvimento, com análise de 
complexidade em cada etapa, além de manutenção de 
invariantes estruturais. A validação será feita tanto 
por demonstrações formais quanto por comparação 
experimental com a reconstrução completa.

\subsubsection{Atualização de atributos dos nós}

A recalculação de atributos como área, momentos e 
caixas delimitadoras pode ser mais custosa do que o 
esperado, sobretudo quando a fragmentação de um nó 
gera múltiplas ramificações. Esse risco pode afetar a 
viabilidade prática do método.

Como contingência, serão identificados atributos que 
podem ser atualizados localmente e 
incrementalmente, e serão propostas estruturas 
auxiliares para evitar revisões globais da 
árvore. Caso seja necessário, a abordagem será 
restrita a um conjunto de atributos e operações de 
menor complexidade.

\subsubsection{Implementação e reprodutibilidade}

A implementação do algoritmo pode revelar problemas 
não observados na formulação matemática, 
como erros de representação da árvore, 
tratamento de casos degenerate e inconsistências 
na manutenção das componentes afetadas.

Como contingência, a implementação será modularizada 
e acompanhada de testes unitários para 
subrotinas específicas. Além disso, serão definidos 
critérios de reprodutibilidade, 
métricas padronizadas e conjuntos de dados de 
referência para avaliar a consistência da solução.

\subsubsection{Validação experimental}

Há risco de que os ganhos observados em termos de 
desempenho computacional não sejam suficientemente 
relevantes em alguns cenários de aplicação, 
reduzindo o impacto prático da proposta.

Como contingência, a avaliação 
experimental será planejada antes da 
implementação, com métricas bem definidas para 
tempo de execução, uso de memória e número de 
operações locais. A comparação com a reconstrução 
completa da \textit{max tree} será central para 
demonstrar o benefício real do operador.

\subsubsection{Cronograma e escopo da pesquisa}

Devido à combinação de teoria, algoritmo, implementação e 
experimentação, o projeto pode sofrer desvio de 
cronograma ou expansão de escopo em uma de suas etapas.

Como contingência, o cronograma será revisado 
periodicamente, e os objetivos serão organizados em 
marcos intermediários. Caso uma extensão teórica ou 
experimental se mostre mais complexa do que o previsto, 
a pesquisa será priorizada para os elementos essenciais 
da contribuição: formalização, correção, eficiência e validação 
do operador de \textit{split}.

\subsubsection{Dependência da literatura}

A literatura sobre atualização dinâmica de 
\textit{component trees} e sobre o operador dual de 
\textit{split} ainda é relativamente escassa, 
o que pode dificultar a comparação direta com o 
estado da arte.

Como contingência, a revisão bibliográfica será contínua e 
integrada ao desenvolvimento do trabalho. 
Sempre que necessário, novas referências serão incorporadas 
ao arcabouço teórico para situar a contribuição no 
contexto da morfologia matemática e das estruturas 
hierárquicas de imagens.